\section{Supplementary Corollary S2.5: Finite-Sample Held-Out Certificate}

S2.4 gives the population lower bound
\[
\operatorname{Cov}(U,S)
\ge
\operatorname{Cov}(Y,S)
-
\sqrt{\mathbb E[(U-Y)^2]\operatorname{Var}(S)}.
\]
This section converts that result into a high-probability certificate from independent held-out data.

Let
\[
(U_i,Y_i,S_i)_{i=1}^n
\]
be i.i.d. held-out observations. The predictor and accessibility rule are assumed fixed before this evaluation sample is drawn, or equivalently the result is applied conditionally on an independently trained fixed model. Assume
\[
|Y|\le B_Y,
\qquad
0\le S\le B_S,
\qquad
|U-Y|\le B_R,
\]
with $B_S>0$ and $0<\mathbb E[S]$.

For $0<\delta<1$, define
\[
\tau_{n,\delta}
=
\sqrt{\frac{\log(10/\delta)}{2n}}.
\]
Let
\[
\bar Y=\frac1n\sum_iY_i,
\qquad
\bar S=\frac1n\sum_iS_i,
\]
\[
\overline{YS}=\frac1n\sum_iY_iS_i,
\qquad
\overline{S^2}=\frac1n\sum_iS_i^2,
\]
and
\[
\widehat M
=
\frac1n\sum_i(U_i-Y_i)^2.
\]
Define
\[
\widehat C
=
\overline{YS}-\bar Y\bar S,
\]
\[
C_L
=
\widehat C-5B_YB_S\tau_{n,\delta},
\]
\[
M_U
=
\min\left\{
B_R^2,
\widehat M+B_R^2\tau_{n,\delta}
\right\},
\]
and
\[
V_U
=
\min\left\{
\frac{B_S^2}{4},
\max\left[
0,
\overline{S^2}
+B_S^2\tau_{n,\delta}
-
\left(\max\{0,\bar S-B_S\tau_{n,\delta}\}\right)^2
\right]
\right\}.
\]
Finally let
\[
D_L=C_L-\sqrt{M_UV_U}.
\]

\begin{corollary}[Finite-sample held-out certificate]
With probability at least $1-\delta$ over the held-out sample,
\[
\operatorname{Cov}(U,S)\ge D_L.
\]
Consequently, if the realized sample satisfies $D_L>0$, then positive population outcome--accessibility covariance is certified at confidence at least $1-\delta$. Under the QBS weighted-measure model, the same event gives
\[
\mathbb E_{\mathrm{FP}}[U]-\mathbb E[U]
\ge
\frac{D_L}{B_S}>0.
\]
\end{corollary}

\begin{proof}
Hoeffding's inequality and the definition of $\tau_{n,\delta}$ imply that each of the following five two-sided deviations fails with probability at most $\delta/5$:
\[
|\mathbb E[Y]-\bar Y|
\le
2B_Y\tau_{n,\delta},
\]
\[
|\mathbb E[S]-\bar S|
\le
B_S\tau_{n,\delta},
\]
\[
|\mathbb E[YS]-\overline{YS}|
\le
2B_YB_S\tau_{n,\delta},
\]
\[
|\mathbb E[S^2]-\overline{S^2}|
\le
B_S^2\tau_{n,\delta},
\]
and
\[
\left|
\mathbb E[(U-Y)^2]-\widehat M
\right|
\le
B_R^2\tau_{n,\delta}.
\]
A union bound gives a simultaneous event of probability at least $1-\delta$.

On this event,
\[
\mathbb E[YS]
\ge
\overline{YS}-2B_YB_S\tau_{n,\delta}.
\]
Furthermore,
\[
\mathbb E[Y]\mathbb E[S]-\bar Y\bar S
=
(\mathbb E[Y]-\bar Y)\mathbb E[S]
+
\bar Y(\mathbb E[S]-\bar S).
\]
Using $\mathbb E[S]\le B_S$ and $|\bar Y|\le B_Y$ gives
\[
\mathbb E[Y]\mathbb E[S]
\le
\bar Y\bar S+3B_YB_S\tau_{n,\delta}.
\]
Thus
\[
\operatorname{Cov}(Y,S)
\ge
C_L.
\]

The squared-residual deviation and the deterministic bound $(U-Y)^2\le B_R^2$ imply
\[
\mathbb E[(U-Y)^2]\le M_U.
\]
For the accessibility variance,
\[
\mathbb E[S^2]
\le
\overline{S^2}+B_S^2\tau_{n,\delta}
\]
and nonnegativity gives
\[
\mathbb E[S]
\ge
\max\{0,\bar S-B_S\tau_{n,\delta}\}.
\]
Therefore the moment expression in $V_U$ is an upper bound on $\operatorname{Var}(S)$. Popoviciu's inequality independently gives
\[
\operatorname{Var}(S)\le\frac{B_S^2}{4},
\]
so
\[
\operatorname{Var}(S)\le V_U.
\]

Applying S2.4 then yields
\[
\operatorname{Cov}(U,S)
\ge
C_L-\sqrt{M_UV_U}
=
D_L.
\]
This proves the population covariance certificate.

If $D_L>0$, T1 gives
\[
\mathbb E_{\mathrm{FP}}[U]-\mathbb E[U]
=
\frac{\operatorname{Cov}(U,S)}{\mathbb E[S]}.
\]
Because $\mathbb E[S]\le B_S$ and the certified covariance lower bound is positive,
\[
\mathbb E_{\mathrm{FP}}[U]-\mathbb E[U]
\ge
\frac{D_L}{B_S}.
\]
\end{proof}

\begin{corollary}[Consistency]
Under the boundedness assumptions above, define the population S2.4 margin
\[
D_*
=
\operatorname{Cov}(Y,S)
-
\sqrt{\mathbb E[(U-Y)^2]\operatorname{Var}(S)}.
\]
If $D_*>0$, then for every fixed $\delta\in(0,1)$,
\[
D_L\to D_*
\]
almost surely as $n\to\infty$. Hence $D_L>0$ eventually almost surely.
\end{corollary}

\begin{proof}
All empirical moments entering $C_L$, $M_U$, and $V_U$ converge almost surely to their population values by the strong law of large numbers, while $\tau_{n,\delta}\to0$. The deterministic caps in $M_U$ and $V_U$ remain valid and converge to quantities no smaller than the corresponding population moments; because the empirical moment-based bounds converge exactly to the relevant population quantities, $D_L\to D_*$. Strict positivity of $D_*$ therefore implies eventual positivity.
\end{proof}

\subsection{Boundary conditions}

The confidence statement requires an independent held-out sample relative to model/accessibility selection. Reusing the same observations for training, model selection, threshold tuning, and certification invalidates the stated guarantee unless that adaptivity is accounted for separately.

The boundedness constants must also be valid before certification. Unbounded or heavy-tailed settings require different concentration tools. Finally, $D_L\le0$ is inconclusive: S2.5 is a sufficient certificate and inherits the conservatism of S2.4 with respect to irreducible conditional outcome variance.
